% SEGMENT OLP-0518-S01
% Part: counterfactuals
% Chapter: introduction
% Section: material-conditional

\documentclass[../../../include/open-logic-section]{subfiles}

\begin{document}

\olfileid{cnt}{int}{mat}

\olsection{பொருள்சார் நிபந்தனைக் கூற்று}

ஆங்கிலத்தில், ஒரு நிபந்தனைக் கூற்றின் மிக எளிய வடிவம் ``If \dots
then \dots'' ஆகும்; இங்கு \dots{} குறிக்கும் பகுதிகளே தனித்தனி
வாக்கியங்கள். எடுத்துக்காட்டாக, ``பணியாளர் அதைச் செய்திருந்தால்,
தோட்டக்காரர் குற்றமற்றவர்.'' அறிமுகத் தருக்கவியல் பாடங்களில்,
நிபந்தனைக் கூற்றுகளை $\lif$ என்ற இணைப்பியால் குறியீடாக்கக் கற்கிறோம்:
\dots{} குறிக்கும் பகுதிகளை, எடுத்துக்காட்டாக !!{formula}s $!A$,~$!B$
ஆகக் குறியீடாக்கி, முழு நிபந்தனைக் கூற்றையும் $!A \lif !B$ என்று
குறியீடாக்குகிறோம்.
% SOURCE CORRECTION TA-CF-001: “earn” is read as “learn.”

இணைப்பி~$\lif$ \emph{மெய்ம்மதிப்புச் சார்ந்தது}; அதாவது $!A \lif
!B$ இன் மெய்ம்மதிப்பு---$\True$ அல்லது $\False$---$!A$,~$!B$
ஆகியவற்றின் மெய்ம்மதிப்புகளால் தீர்மானிக்கப்படுகிறது: $!A$ பொய்யாகவோ
$!B$ மெய்யாகவோ இருந்தால், எனில் மட்டுமே, $!A \lif !B$ மெய்யாகும்;
மற்றபடி அது பொய்யாகும். மெய்ம்மதிப்பு மதிப்பீடு~$\pAssign v$ ஐப்
பொறுத்து, $\pSat/{v}{!A}$ அல்லது $\pSat{v}{!B}$ எனில், எனில் மட்டுமே,
$\pSat{v}{!A \lif !B}$ என்று வரையறுக்கிறோம். இந்தப் பொருண்மையியலைக்
கொண்ட இணைப்பி~$\lif$ \emph{பொருள்சார் நிபந்தனை} எனப்படுகிறது.

இந்த வரையறையிலிருந்து பல அடிப்படைத் தருக்க உண்மைகள் கிடைக்கின்றன.
முதலில், பொருள்சார் நிபந்தனைக்கு வருவித்தல் தேற்றம் பொருந்துகிறது:
\begin{align}
  \text{எனில் } \Gamma, !A \Entails !B \text{ அப்போது } \Gamma \Entails !A \lif !B
\end{align}
இது மெய்ம்மதிப்புச் சார்ந்தது: $!A \lif !B$ உம் $\lnot !A \lor !B$ உம்
சமானமானவை:
\begin{align}
  !A \lif !B & \Entails \lnot !A \lor !B\\
  \lnot !A \lor !B & \Entails !A \lif !B
  \intertext{ஒரு பொருள்சார் நிபந்தனைக் கூற்று அதன் பின்னிலையிலிருந்தும்,
    அதன் முன்னிலையின் மறுப்பிலிருந்தும் பின்விளைகிறது:}
  !B & \Entails !A \lif !B\\
  \lnot !A & \Entails !A \lif !B
  \intertext{பொய்யான ஒரு பொருள்சார் நிபந்தனைக் கூற்று, அதன் முன்னிலையும்
    பின்னிலையின் மறுப்பும் சேர்ந்த இணைப்பிற்கு சமானமானது: $!A \lif !B$
    பொய்யெனில் $!A \land \lnot !B$ மெய்; மறுதிசையிலும் அவ்வாறே:}
  \lnot(!A \lif !B) & \Entails !A \land \lnot !B\\
  !A \land \lnot !B & \Entails \lnot(!A \lif !B)
  \intertext{பொருள்சார் நிபந்தனை மோடஸ் போனென்ஸை ஆதரிக்கிறது:}
  !A, !A \lif !B & \Entails !B
  \intertext{பொருள்சார் நிபந்தனை ஒருங்குவிக்கிறது:}
  !A \lif !B,  !A \lif !C & \Entails !A \lif (!B \land !C)
  \intertext{முன்னிலையை நாம் எப்போதும் வலுப்படுத்தலாம்; அதாவது, இந்த
    நிபந்தனை \emph{ஒருபோக்குத்தன்மையுள்ளது}:}
  !A \lif !B & \Entails (!A \land !C) \lif !B
  \intertext{பொருள்சார் நிபந்தனை கடப்புத்தன்மையுள்ளது; அதாவது, சங்கிலி
    விதி செல்லுபடியாகும்:}
  !A \lif !B, !B \lif !C & \Entails !A \lif !C
  \intertext{பொருள்சார் நிபந்தனை அதன் எதிர்மறைநேர்மாற்றுக் கூற்றுக்குச்
    சமானமானது:}
  !A \lif !B & \Entails \lnot !B \lif \lnot !A\\
  \lnot !B \lif \lnot !A & \Entails !A \lif !B
\end{align}

இவையெல்லாம் கணிதப் பகுத்தறிதலில் பயனுள்ள, சிக்கலற்ற அனுமானங்கள்.
ஆனால், கணிதமல்லாத சூழல்களில் இதற்கு ஒத்த அனுமானங்களுக்கு எதிரெடுத்துக்காட்டுகள்
என்று கூறப்படுபவை மெய்யியல் மற்றும் மொழியியல் இலக்கியங்களில் நிறைந்துள்ளன.
ஆங்கில ``if \dots then \dots'' கூற்றுகளைக் குறியீடாக்கும்போது பயன்படுத்த
வேண்டிய பொருத்தமான இணைப்பி பொருள்சார் நிபந்தனை~$\lif$ அல்ல---அல்லது
குறைந்தது எல்லா நேரங்களிலும் அல்ல---என்பதை அவை சுட்டுகின்றன.

\end{document}
